% On-policy distillation (OPD) queries a teacher on responses sampled from the student itself \citep{agarwal2024onpolicy}. Unlike a fixed collection of teacher targets, this supervision changes with the student's predictions and the prefixes its rollouts visit. Knowledge and sequence distillation establish the transfer setting \citep{hinton2015distilling,kim2016sequence}; divergence objectives and response sources further shape how feedback is used \citep{gu2023minillm,ko2024distillm,ko2025distillm2}. We call the teacher-induced loss and gradients at the current training state \emph{residual supervision}. A falling Kullback--Leibler (KL) loss leaves a substantive question: what supervision remains for the student to learn from? Its organization matters because energy concentrated on a few tokens can still point along many hidden directions, and examples varying along common axes can request different signed corrections. Figure~\ref{fig:residual_overview} introduces these distinctions.

% \begin{figure*}[t]
% \centering
% \includegraphics[width=\textwidth]{figures/residual_overview.pdf}
% \caption{Residual supervision exhibits distinct forms of concentration and agreement. The panels summarize paired FULL measurements at both scales: endpoint token-energy concentration; endpoint displacement and gradient coverage by the same fixed early rank-128 basis; and early/late top-4 covariance-axis overlap versus mean-gradient cosine. Endpoint measurements use update 106; agreement compares steps 10 and 100. Full trajectories and intermediate checkpoints appear in Figures~\ref{fig:two_scale_dynamics}--\ref{fig:two_scale_batch}. Historical displacement and current gradients use different prefix populations.}
% \label{fig:residual_overview}
% \end{figure*}

% Existing analyses address complementary objects: data coverage and supervision absorption \citep{li2026rethinking,fu2026rethinking}, compact adaptation and parameter-update geometry \citep{li2018intrinsic,aghajanyan2021intrinsic,hu2021lora,shen2026geometry,yu2026dense}, and token importance and feedback reliability \citep{xu2026tip,sheng2026ier,zhao2026vanishing,jiang2026rock,armandpour2026unmasking}. We examine the joint organization of the \emph{current} teacher-induced gradients along matched OPD trajectories. Historical representation change supplies a reference for current sensitivity, while signed gradient agreement tests whether different examples request similar average corrections, a distinction motivated by work on gradient confusion and coherence \citep{sankararaman2020confusion,chatterjee2020coherent}.

% On mathematical reasoning with Qwen3-1.7B and Qwen3-4B students, strong disagreements become rarer while the current highest-KL decile retains most hidden-gradient energy (Figure~\ref{fig:two_scale_dynamics}). Concentration over positions coexists with current sensitivity extending beyond historical adaptation: the same early rank-128 basis covers most accumulated representation displacement but a much smaller share of current gradient energy (Figure~\ref{fig:subspace}; Tables~\ref{tab:paired_adapter_geometry} and~\ref{tab:paired4b_adapter_geometry}). Across training subsets, dominant axes of gradient variation remain shared as signed mean-gradient agreement weakens (Figure~\ref{fig:two_scale_batch}). Coherence within the early basis also loses much of its initial relative advantage (Table~\ref{tab:space_coherence}). Together, these contrasts show that residual supervision can remain concentrated and structured while examples become less aligned in their average corrections.

% We extend this picture to executed optimizer steps and practical resource choices. Holding validation responses fixed across each before/after measurement reveals diminishing held-out teacher-KL gains; parameter-step diagnostics track accompanying changes in magnitude, alignment, and finite-step residuals (Figure~\ref{fig:two_scale_transfer}). These measurements characterize the actual update effects alongside the local gradients. Finally, stopping, compact continuation, and response reuse provide exploratory performance--cost comparisons (Table~\ref{tab:benchmark_summary}, Figure~\ref{fig:efficiency_overview}), drawing on compact adaptation and cached-data training \citep{hu2021lora,kang2026async}. Their distinct resource savings establish practical context for the empirical study of remaining teacher supervision.
\section{Introduction}
\begin{description}[leftmargin=1.2em,style=nextline]
\item[\textbf{RQ1: What drives the learning signal, and how does it change during training?}]
\item[\textbf{RQ2: From a representation perspective, where does the learning signal settle?}]
\item[\textbf{RQ3: How much of the signal transfers to unseen responses?}]
\item[\textbf{RQ4: Can OPD training be made more efficient?}]
\end{description}